\begin{textsample}{NEG Dim 3 – go – Score -5.00 – go/go\_em\_42.txt}  \label{ex:f3_neg_242}
1.3 - Estrutura do Quadro de \textbf{Competências} e Habilidades Os eixos orientadores da estruturação, concepção e organização de todos os Objetivos de Aprendizagem aqui apresentados foram construídos a partir das \textbf{competências} e habilidades específicas que compõem a área de Linguagens e suas Tecnologias na Base Nacional Comum Curricular ( BNCC ) referentes à Etapa do Ensino Médio.
O quadro de \textbf{competências} detalha e referencia os caminhos didático-pedagógicos para o alcance das habilidades, tendo como parâmetro o sistema organizacional para o Ensino Médio, proposto pela BNCC ( BRASIL, 2018 ), optou -se pela distinção, no interior da Área de Linguagens e suas Tecnologias, dos cinco componentes curriculares : Arte, Educação Física e Línguas \textbf{espanhola}, Inglesa e Portuguesa - tendo em vista que determinadas \textbf{competências} tendem a contemplar alguns componentes em específico.
Desse modo, a sequência de apresentação das \textbf{competências} e habilidades se estruturam : - \textbf{Competências} 1, 2 e 3 - Língua Portuguesa, Arte, Educação Física, Línguas \textbf{Espanhola} e Inglesa - \textbf{Competência} 4 - Línguas \textbf{Espanhola} e Inglesa - \textbf{Competência} 5 - Educação Física - \textbf{Competência} 6 - Arte - \textbf{Competência} 7 - TDICs As tabelas foram subdivididas em cinco colunas, nomeadas, da esquerda para a direita, da seguinte forma : HABILIDADES, OBJETIVOS DE APRENDIZAGEM, CAMPOS DE ATUAÇÃO, PRÁTICAS DE LINGUAGEM e OBJETOS DE CONHECIMENTO.
Segue recorte do cabeçalho.
O quadro 06 traz a primeira \textbf{competência} para ilustrar como a área de Linguagens e suas Tecnologias estruturou o quadro de \textbf{competências} e habilidades : Os cincos componentes relacionam -se com os campos de atuação social, e estão vinculados entre si por meio das habilidades, das práticas de linguagens e dos objetos de conhecimento, uma vez que, na exploração expressiva das diferentes linguagens, o que se busca é uma abordagem multissemiótica, o que inclui elementos visuais, sonoros, verbais e corporais.
Na primeira coluna, reuniram -se as HABILIDADES apresentadas na BNCC ( Brasil, 2018 ), por isso estão mencionadas como “ Habilidades da BNCC ”.
Em nossa grande área é possível encontrar habilidade associada a Língua Portuguesa e identificadas com a sigla \textbf{LP}.
Na coluna OBJETIVOS DE APRENDIZAGEM ( OA ), a equipe redatora elaborou um conjunto de capacidades ou de aplicações, obedecendo à sequência dos sete processos cognitivos das Taxonomia de Bloom e relacionando -os às habilidades da coluna anterior.
São conhecimentos e qualificações a serem atingidos pelos/as estudantes, com o propósito de alcançar, ao final, as respectivas \textbf{competências} em sua plenitude.
Segue um exemplo de Objetivo de Aprendizagem, como disposto no Documento Curricular para Goiás - Etapa Ensino Médio :

% matched lemmas: competência, espanhol, lp
\end{textsample}
